% =====================================================================
%  CSE 200 Online-1 (C2)  ---  "The Solar System"
%  Faithful reproduction. Compile with pdflatex TWICE (TOC + refs).
%  Replace the example-image-* placeholders with the real planet images
%  supplied in the exam (sun.jpg, mercury.jpg, ...).
% =====================================================================
\documentclass[10pt]{article}

% ---------- Geometry / encoding ----------
\usepackage[a4paper,margin=1in]{geometry}
\usepackage[T1]{fontenc}
\usepackage[utf8]{inputenc}

% ---------- Mathematics ----------
\usepackage{amsmath}   % align, \tag, \implies, \dfrac, \text
\usepackage{amssymb}   % \odot, \varpi, \diamond, \dagger, \propto
\usepackage{amsthm}    % proposition + proof environments

% ---------- Figures & subfigures ----------
\usepackage{graphicx}
\usepackage{subcaption}    % subfigure environment + \subref  (1a,1b,...)

% ---------- Colour tables ----------
\usepackage[table,dvipsnames]{xcolor} % loads colortbl; \rowcolor/\cellcolor
\usepackage{array}
\usepackage{multirow}

% ---------- Lists ----------
\usepackage{enumitem}      % (\Roman*), $\dagger$, $\diamond$ labels

% ---------- Author affiliation ----------
\usepackage{authblk}

% ---------- Hyperlinks (LAST) ----------
\usepackage{hyperref}
\hypersetup{
  colorlinks = true,
  linkcolor  = blue,   % clickable TOC / \ref
  citecolor  = blue,
  urlcolor   = magenta
}

% ---------- Theorem style ----------
\newtheorem{proposition}{Proposition}[section]  % numbers like 3.1

% =====================================================================
\title{The Solar System}
\author[1]{Johannes Kepler}
\affil[1]{Imperial Mathematician, Prague}
\date{}
% =====================================================================

\begin{document}
\maketitle
\tableofcontents
\bigskip

% =====================================================================
\section{Our Planetary Neighborhood}

\subsection{The Sun and Its Planets}
The Sun and its gravity-bound bodies form the Solar System~\cite{kepler}.
See \href{https://solarsystem.nasa.gov/}{NASA Solar System Exploration}.
In increasing distance, the planets are
\[
  \text{Mercury, Venus, Earth, Mars, Jupiter, Saturn, Uranus, Neptune.}
\]
The planets form two groups:
\begin{enumerate}[label=(\Roman*)]
  \item \textbf{Inner planets}
    \begin{itemize}[label=$\dagger$]
      \item Mercury, Venus, Earth, and Mars belong to this group.
      \item They are comparatively small and have solid surfaces.
      \item They are also called the terrestrial planets.
    \end{itemize}
  \item \textbf{Outer planets}
    \begin{itemize}[label=$\dagger$]
      \item Jupiter and Saturn are gas giants.
      \item Uranus and Neptune are commonly called ice giants.
      \item All four outer planets possess ring systems.
    \end{itemize}
\end{enumerate}
The inner planets are small, dense, and relatively close to the Sun,
whereas the outer planets are much larger and travel along considerably
wider orbits.

\subsection{Portraits of the Solar System}
Figure~\ref{fig:solar} shows the Sun and eight planets.\footnote{Sizes and
distances are schematic.}
Compare Mars in Figure~\ref{fig:mars} with Jupiter in Figure~\ref{fig:jupiter}.

\begin{figure}[htbp]
  \centering

  {\large\textbf{Star}}\par\medskip
  \begin{subfigure}[b]{0.30\textwidth}
    \centering
    \includegraphics[width=0.6\linewidth]{example-image-a}
    \caption{The Sun}\label{fig:sun}
  \end{subfigure}

  \bigskip
  {\large\textbf{Inner Planets}}\par\medskip
  \begin{subfigure}[b]{0.23\textwidth}
    \includegraphics[width=\linewidth]{example-image-b}
    \caption{Mercury}\label{fig:mercury}
  \end{subfigure}\hfill
  \begin{subfigure}[b]{0.23\textwidth}
    \includegraphics[width=\linewidth]{example-image-c}
    \caption{Venus}\label{fig:venus}
  \end{subfigure}\hfill
  \begin{subfigure}[b]{0.23\textwidth}
    \includegraphics[width=\linewidth]{example-image-a}
    \caption{Earth}\label{fig:earth}
  \end{subfigure}\hfill
  \begin{subfigure}[b]{0.23\textwidth}
    \includegraphics[width=\linewidth]{example-image-b}
    \caption{Mars}\label{fig:mars}
  \end{subfigure}

  \bigskip
  {\large\textbf{Outer Planets}}\par\medskip
  \begin{subfigure}[b]{0.23\textwidth}
    \includegraphics[width=\linewidth]{example-image-c}
    \caption{Jupiter}\label{fig:jupiter}
  \end{subfigure}\hfill
  \begin{subfigure}[b]{0.23\textwidth}
    \includegraphics[width=\linewidth]{example-image-a}
    \caption{Saturn}\label{fig:saturn}
  \end{subfigure}\hfill
  \begin{subfigure}[b]{0.23\textwidth}
    \includegraphics[width=\linewidth]{example-image-b}
    \caption{Uranus}\label{fig:uranus}
  \end{subfigure}\hfill
  \begin{subfigure}[b]{0.23\textwidth}
    \includegraphics[width=\linewidth]{example-image-c}
    \caption{Neptune}\label{fig:neptune}
  \end{subfigure}

  \caption{The Sun and the eight planets of the Solar System, arranged
  into the inner and outer planetary groups. The displayed sizes and
  distances are not to scale.}
  \label{fig:solar}
\end{figure}

% =====================================================================
\section{Kepler's Laws of Planetary Motion}
Kepler's three laws describe planetary motion~\cite{nasa-kepler}.

\subsection{First Law: Elliptical Orbits}
A planetary orbit is elliptical, with the Sun at one focus:
\begin{equation}
  \frac{x^{2}}{a^{2}} + \frac{y^{2}}{b^{2}} = 1,
  \qquad a \ge b > 0 .
\end{equation}
Here $a$ and $b$ are the semi-axes.

For focal distance $c$,
\begin{equation}
  \begin{aligned}
    c^{2} &= a^{2} - b^{2} \\
    e     &= \frac{c}{a}
  \end{aligned}
\end{equation}
Here $e$ is the eccentricity.

In polar form,
\begin{equation}
  r(\theta) = \frac{a\,(1 - e^{2})}{1 + e\cos\theta}.
  \label{eq:polar}
\end{equation}
Equation~\eqref{eq:polar} gives the orbital distance.

\subsection{Second Law: Equal Areas in Equal Times}
Equal times sweep equal areas:
\begin{equation}
  \Delta t_{1} = \Delta t_{2}
  \implies
  \Delta A_{1} = \Delta A_{2}.
\end{equation}
Consequently, a planet moves
\begin{itemize}[label=$\diamond$]
  \item faster when it is closer to the Sun, and
  \item slower when it is farther from the Sun.
\end{itemize}

\subsection{Third Law: Orbital Period and Distance}
For planets orbiting one star,
\begin{equation}
  T^{2} \propto a^{3}.
\end{equation}
Newtonian form:
\[
  T^{2} = \frac{4\pi^{2}}{G M_{\odot}}\, a^{3},
\]
Here $M_{\odot}$ is the solar mass.

\subsection{An Earth--Mars Comparison}
Using Earth as reference,
\[
  a_{E} = 1~\text{AU}, \qquad T_{E} = 1~\text{year}
\]
where $1~\text{AU} \approx 1.496 \times 10^{11}~\text{m}$.\footnote{AU
denotes the mean Earth--Sun distance.} For Mars, $a_{M} \approx
1.524~\text{AU}$, so
\[
  \frac{T_{M}^{2}}{T_{E}^{2}} = \frac{a_{M}^{3}}{a_{E}^{3}},
  \qquad
  T_{M} = T_{E}\left(\frac{a_{M}}{a_{E}}\right)^{3/2}
        \approx (1.524)^{3/2} \approx 1.88~\text{years}.
\]
Thus $T_{M} \approx 1.88$ years.

% =====================================================================
\section{A Relativistic Extension}
General relativity extends Kepler's model.\footnote{Planetary corrections
are small but measurable.} Einstein's equations are
\begin{equation}
  G_{\mu\nu} + \Lambda g_{\mu\nu} = \frac{8\pi G}{c^{4}}\, T_{\mu\nu}.
\end{equation}
For spherical mass $M$,
\begin{equation}
  ds^{2} = -\left(1 - \frac{r_{s}}{r}\right)c^{2}\,dt^{2}
           + \left(1 - \frac{r_{s}}{r}\right)^{-1} dr^{2}
           + r^{2}\left(d\theta^{2} + \sin^{2}\!\theta\, d\varphi^{2}\right),
  \qquad r_{s} = \frac{2GM}{c^{2}}.
\end{equation}
Free fall satisfies
\begin{equation}
  \frac{d^{2}x^{\mu}}{d\tau^{2}}
  + \Gamma^{\mu}_{\alpha\beta}\,
    \frac{dx^{\alpha}}{d\tau}\frac{dx^{\beta}}{d\tau} = 0 .
\end{equation}
Per orbit,
\begin{equation}
  \Delta\varpi = \frac{6\pi G M}{a\,(1 - e^{2})\,c^{2}} .
\end{equation}

\begin{proposition}[Weak-field circular orbit]\label{prop:weak}
For a circular orbit of radius $r$,
\begin{equation}
  \left(\frac{v}{c}\right)^{2} = \frac{GM}{r c^{2}} = \frac{r_{s}}{2r}.
  \tag{10}
\end{equation}
\end{proposition}

\begin{proof}
The circular-orbit balance $v^{2}/r = GM/r^{2}$ gives $v^{2} = GM/r$.
Divide by $c^{2}$ and use $r_{s} = 2GM/c^{2}$.
\end{proof}

See Proposition~\ref{prop:weak}.

% =====================================================================
\section{The Outer Planets}
The outer planets are Jupiter, Saturn, Uranus, and
Neptune~\cite{nasa-planets}. The first two are gas giants; the last two
are ice giants.

Table~\ref{tab:outer} uses scientific and relativistic
notation.\footnote{Values are rounded for typesetting practice.}

\begin{table}[htbp]
  \centering
  \caption{Orbital and relativistic parameters of the outer planets}
  \label{tab:outer}
  \renewcommand{\arraystretch}{1.5}
  \begin{tabular}{|l|cc|cc|}
    \hline
    \rowcolor{Gray!35}
    \multirow{2}{*}{\textbf{Planet}}
      & \multicolumn{2}{c|}{\textbf{Orbital parameters}}
      & \multicolumn{2}{c|}{\textbf{Relativistic notation}} \\
    \cline{2-5}
    \rowcolor{Gray!20}
      & $a$ (AU) & $e$
      & $\varepsilon = \dfrac{1}{a c^{2}} G M_{\odot}$
      & $\beta = \sqrt{\varepsilon}$ \\
    \hline
    \rowcolor{Yellow!15} Jupiter & 5.204  & $4.89\times10^{-2}$ & $1.90\times10^{-9}$  & $4.36\times10^{-5}$ \\
    \rowcolor{Yellow!25} Saturn  & 9.583  & $5.65\times10^{-2}$ & $1.03\times10^{-9}$  & $3.21\times10^{-5}$ \\
    \rowcolor{Cyan!15}   Uranus  & 19.191 & $4.72\times10^{-2}$ & $5.14\times10^{-10}$ & $2.27\times10^{-5}$ \\
    \rowcolor{Cyan!25}   Neptune & 30.070 & $8.6\times10^{-3}$  & $3.28\times10^{-10}$ & $1.81\times10^{-5}$ \\
    \hline
  \end{tabular}
\end{table}

Table~\ref{tab:outer} compares the four giant planets.

% =====================================================================
\begin{thebibliography}{9}
  \bibitem{kepler} J.~Kepler, \textit{Astronomia Nova}. Heidelberg, 1609.
  \bibitem{nasa-kepler} National Aeronautics and Space Administration,
    ``Orbits and kepler's laws,'' 2024.
  \bibitem{nasa-planets} National Aeronautics and Space Administration,
    ``About the planets,'' 2026.
    \href{https://solarsystem.nasa.gov/}{Solar System Exploration}.
\end{thebibliography}

\end{document}
